% Limitations (task S3, Person A). Items about RQ1 and RQ2 belong to Person B
% and go in the marked place before this goes into the paper.

\section{Limitations}
\label{sec:limits}
\begin{itemize}
  \item \emph{Few stores.} The retail clips come from a small number of cameras
        in a small number of stores. The intervals in Section~\ref{sec:rq3} cover
        the variation between people, not between stores, and a store with
        different lighting or camera placement may behave differently.
  \item \emph{Labels are judgements.} The ground truth is what two annotators
        saw in the crops. People they marked unsure are left out, which removes
        the hardest cases from every accuracy we report.
  \item \emph{Correlated frames.} Crops of one person taken 0.4\,s apart share
        lighting, pose and clothing, so their errors are not independent. Pooling
        cannot fix an error that every frame makes, and accuracy against the
        number of frames pooled shows how much independent evidence the frames
        really hold.
  \item \emph{Hand-chosen cues.} The quality model sees seven cues that we chose.
        A failure that none of them measures, such as a bag held in front of the
        body or strong backlight, gets no lower weight.
  % If the MARS images are obtained and used (--attention-train), rewrite this item.
  \item \emph{Small training sets.} With MARS unavailable, the quality model and
        the attention head were fitted on the retail people themselves, with
        cross-validation. Both would likely improve with more data, and the
        attention head, with far more parameters, more so.
  \item \emph{The tracker.} Pooling is only as good as the track it pools. The
        tracker has no appearance model, so two people who cross paths can swap
        tracks, and a pooled label can then mix two people.
  % Person B: add the RQ1 and RQ2 limitations here (PETA split, single store as
  % FDA target, seeds).
\end{itemize}
